\concept{B01}{Sensor samples and a sampling clock}{F04, F08, F14}
\paragraph{Foundation and definition.}
A continuous physical signal and its recorded samples are different objects.
A sampling model chooses discrete times and records a finite vector at each
time. For the six-channel inertial input, write
\[
 v[n]=(a_x[n],a_y[n],a_z[n],\omega_x[n],\omega_y[n],\omega_z[n])^\top,
 \qquad t_n=t_0+n/f_s.
\]
\paragraph{Symbols and meaning.}
$n\in\{0,\ldots,T-1\}$ is a sample index, $T$ is the number of samples,
$v[n]\in\R^6$ is one observation, $a_x,a_y,a_z$ are accelerometer axes,
and $\omega_x,\omega_y,\omega_z$ are angular-velocity axes. $t_0$ is the
first sample time in seconds and $f_s>0$ is samples per second. HAPT's
nominal $f_s=50$ corresponds to $1/f_s=0.02$ seconds per sample. Its raw
acceleration unit is g; angular velocity uses rad/s \cite{hapt}.
\paragraph{Derivation and example.}
For $T=128$, the first-to-last elapsed interval is
$t_{127}-t_0=127/50=2.54$ s. The sample-count convention assigns
$T/f_s=2.56$ s. These are different measurements, not contradictory
estimates. Adding a sampling interval to the last timestamp is a convention,
not another observed sample. Raw recordings without timestamps do not
independently establish a perfectly uniform clock.
\paragraph{Why ECDR needs it.}
A query stride measured in samples becomes a physical query rate only after
choosing this clock model. A receptive field of nine samples is not nine
seconds. Equal sample identity is also the foundation of exact overlap.
\paragraph{Conditions and connections.}
Uniform sampling is a modeling assumption here, not a consequence of the
field axioms. Timing jitter, gaps and asynchronous sensors require a richer
input contract. B02 defines temporal span; B04 defines the query grid.

\concept{B02}{Temporal span: what an annotated interval covers}{F08, F24, B01}
\paragraph{Foundation and definition.}
A temporal annotation span is an interval of sample indices assigned one
activity. It is not the linear span of vectors in F12. With inclusive
endpoints $a,b\in\mathbb Z$ and $a\le b$, the index set is
\[
 I_{a,b}=\{n\in\mathbb Z:a\le n\le b\},\qquad |I_{a,b}|=b-a+1.
\]
\paragraph{Symbols and meaning.}
$a$ and $b$ denote the first and last annotated sample indices; $n$ ranges
over samples; $I_{a,b}$ is a finite set; $|I_{a,b}|$ is its cardinality,
not a norm. With sampling rate $f_s>0$, sample-count duration is
$(b-a+1)/f_s$ seconds; first-to-last elapsed time is $(b-a)/f_s$ seconds.
The half-open representation $[a,b+1)$ denotes the same integer samples.
\paragraph{Derivation and example.}
The span $[64,191]$ contains $191-64+1=128$ samples. At 50 Hz it has
sample-count duration $2.56$ s and elapsed time $2.54$ s. A span $[1,128]$
has the same length, but a window grid with starts $0,64,128,\ldots$ does
not contain a length-128 window inside it. Duration and grid alignment
are separate constraints, derived in B05.
\paragraph{Why ECDR needs it.}
Annotation support after windowing depends on which windows fit inside a
span. Short transitions may disappear from a pure-window evaluation even
when their annotations remain in the raw data. Understanding the unit of
labeling prevents confusing a class transition with a mixed-label window.
\paragraph{Conditions and connections.}
Index origin and endpoint convention must be verified from source metadata.
The inherited HAPT pipeline reproduces one convention; reproduction does
not independently prove its original annotation semantics. B21 uses the
word span for a receptive-field extent, another distinct meaning.

\concept{B03}{A tensor, vectorization and a dimension check}{F14, F15, F19}
\paragraph{Foundation and definition.}
A time-series tensor organizes values by named axes. For one window,
$X\in\R^{C\times W}$ has $C$ channels and $W$ sample positions. A
channel-major vectorization is the bijection
\[
 x_{cW+i}=X[c,i],\qquad
 c=0,\ldots,C-1,\quad i=0,\ldots,W-1.
\]
\paragraph{Symbols and meaning.}
$X[c,i]$ is one scalar observation; $x=\operatorname{vec}(X)\in\R^d$ is
the same collection in a vector; $d=CW$ is input dimension. $c$ indexes
channels and $i$ indexes within-window time. A batch adds an axis:
$X^{\mathrm{batch}}\in\R^{B\times C\times W}$ with $B$ windows. Batch
index is not a new physical sensor or a new time sample.
\paragraph{Derivation and example.}
For $C=2,W=3$, rows $(1,2,3)$ and $(4,5,6)$ become
$x=(1,2,3,4,5,6)^\top$. Another fixed layout is a permutation $Px$,
where $P\in\R^{d\times d}$ has one 1 per row and column. Since
$P^\top P=I_d$, $\norm{Px}_2=\norm{x}_2$. Here $I_d$ is the identity
matrix. However, unchanged CNN weights generally do not compensate for an
arbitrary input permutation.
\paragraph{Why ECDR needs it.}
The classifier, domain generator, normalization and membership checker must
agree on the layout. For HAPT, $C=6,W=128$ gives $d=768$; the final feature
dimension may be 64, a different object. A float32 input has
$768\times4=3072$ payload bytes, not total runtime memory.
\paragraph{Conditions and connections.}
Vectorization is lossless only when its inverse and axis convention are
fixed. Unit consistency is addressed by B07--B08; actual CNN operators
consume the structured tensor rather than an unlabeled coordinate list.

\concept{B04}{The sliding-window grid and its exact count}{F08, F09, B01}
\paragraph{Foundation and definition.}
A no-padding query grid takes complete windows from one recording. For
$T$ samples, window length $W\ge1$ and stride $s\ge1$, starts are $ks$,
where $k$ is a nonnegative integer. A valid window reads
\[
 X_k[:,i]=v[ks+i],\qquad i=0,\ldots,W-1,
 \qquad ks+W\le T.
\]
\paragraph{Symbols and meaning.}
$v[n]\in\R^C$ is sample $n$, $C$ is channel count, $X_k\in\R^{C\times W}$
is query $k$, and ``:'' means every channel. $T,W,s,k,i$ count samples
or windows, not seconds. No window crosses a recording boundary.
\paragraph{Derivation and example.}
The final inequality is equivalent to $k\le(T-W)/s$. When $T\ge W$,
\[
 N_T=1+\left\lfloor\frac{T-W}{s}\right\rfloor;
 \qquad N_T=0\text{ if }T<W.
\]
There are $N_T$ integers because index 0 is included. At
$T=1000,W=128,s=64$, $N_T=14$. The last start is $13\cdot64=832$,
ending at exclusive index 960; the next would end at 1024 and is invalid.
The nominal query rate is $f_s/s=50/64=0.78125$ queries/s within a long
uniformly sampled recording, ignoring startup and termination.
\paragraph{Why ECDR needs it.}
Query count sets evaluation denominators and amortization frequency.
Overlapping windows do not multiply the independent recording duration.
For multiple recordings, sum their individual counts; there is no valid
pair from the end of one recording to the beginning of another.
\paragraph{Conditions and connections.}
This count assumes fixed length, fixed stride, no padding and a recording
grid starting at index 0. Different alignment policies require a different
count. B05 applies the same integer reasoning to an annotation span.

\concept{B05}{How many grid windows fit inside a label span?}{F08, B02, B04}
\paragraph{Foundation and definition.}
A pure window must fit wholly inside one consistently labeled interval.
Use a half-open span $[a,b)$ inside a recording with indices $[0,T)$.
The grid starts remain $ks$; they are not reset at each annotation.
\paragraph{Symbols and meaning.}
$a,b,T\in\mathbb Z$ are span and recording bounds, $W\ge1$ is window
length, $s\ge1$ is stride, and $k\ge0$ is a window index. Define integer
bounds $k_{\min}$ and $k_{\max}$ by
\[
 k_{\min}=\left\lceil\frac{\max(a,0)}s\right\rceil,\qquad
 k_{\max}=\left\lfloor\frac{\min(b-W,T-W)}s\right\rfloor.
\]
\paragraph{Derivation and example.}
The required inequalities are $ks\ge a$, $ks+W\le b$, $ks\ge0$ and
$ks+W\le T$. Dividing by positive $s$, then rounding lower bounds upward
and upper bounds downward, yields
\[
 N_{[a,b)}=\max(0,k_{\max}-k_{\min}+1).
\]
At $a=100,b=400,W=128,s=64,T\ge400$, the bounds are 2 and 4, so starts
128,192,256 are valid. Start 320 ends at 448 and is outside the span.
For inclusive input $[a,b_{\rm inc}]$, substitute $b=b_{\rm inc}+1$.
\paragraph{Why ECDR needs it.}
A span with at least $W$ samples can still produce zero pure windows if
the fixed grid misses it. Thus pure class support is a property of both
the annotated signal and the query policy. It does not directly estimate
the prevalence of that activity in the raw recording or in deployment.
\paragraph{Conditions and connections.}
The formula treats one span and one recording. Adjacent spans with the
same activity require an explicit merge policy before adding counts.
Changing the index convention without validation changes the population
being analyzed; an algebraically correct formula cannot validate that choice.

\concept{B06}{A label policy is a definition of the evaluation population}{F03, F08, B05}
\paragraph{Foundation and definition.}
Window labels summarize many sample-level labels. For $K$ activity classes,
let $n_j$ count samples assigned class $j\in\{1,\ldots,K\}$ and let $n_0$
count unannotated samples. Then
\[
 \sum_{j=0}^{K}n_j=W,\qquad
 p_{\max}=\max_{1\le j\le K}n_j/W.
\]
\paragraph{Symbols and meaning.}
$W$ is the number of samples in the window; each $n_j$ is an integer
between 0 and $W$. $p_{\max}$ is a whole-window proportion, not a CNN
probability. ID0 denotes absent annotation here, not an additional
recognized activity. An inherited policy is pure if some $n_j=W$;
unlabeled if $n_0=W$; otherwise majority when $p_{\max}\ge0.6$, and
ambiguous when $p_{\max}<0.6$.
\paragraph{Derivation and example.}
Because counts are integer, a majority threshold of 0.6 at $W=128$ needs
$n_j\ge\lceil0.6\cdot128\rceil=77$. If 70 samples have one activity and
58 are unannotated, the observed-label proportion is $70/70=1$, but the
policy proportion is $70/128=0.546875$: the window is ambiguous. The
choice of denominator changes the definition, not just its display.
\paragraph{Why ECDR needs it.}
Filtering to pure windows changes which transitions and mixed boundaries
remain in a study. A reference can be trained on a subset while a runtime
stream still contains every query. Reuse coverage needs its own explicitly
stated population. True activity labels are available for retrospective
analysis, not automatically to a runtime event selector.
\paragraph{Conditions and connections.}
Tie handling for a majority ID and annotation conventions must be fixed.
The policy is a modeling choice; probability axioms do not prescribe it.
B30 and C28 explain why metrics with different denominators cannot be renamed.

\concept{B07}{Frozen normalization as an invertible affine map}{F15--F17, B03}
\paragraph{Foundation and definition.}
Normalization changes coordinates rather than changing the underlying
recorded observation. In vector form, fixed per-coordinate centering and
positive scaling give
\[
 z=\mathcal P(x)=D_\sigma^{-1}(x-\mu),\qquad
 D_\sigma=\operatorname{diag}(\sigma_1,\ldots,\sigma_d).
\]
\paragraph{Symbols and meaning.}
$x\in\R^d$ is raw input, $z\in\R^d$ is normalized input,
$\mu\in\R^d$ contains fitted means, and every $\sigma_i>0$ is a fixed
scale. $D_\sigma\in\R^{d\times d}$ is diagonal. For per-channel fitting,
the same mean and scale repeat across positions of that channel.
\paragraph{Derivation and example.}
Invertibility follows directly: $x=D_\sigma z+\mu$. If one coordinate has
$\mu_i=2,\sigma_i=0.5$ and raw value $x_i=2.1$, then $z_i=0.2$.
For anchor $x^r$ and its image $z^r=\mathcal P(x^r)$,
\[
 z-z^r=D_\sigma^{-1}(x-x^r).
\]
Centering cancels in a difference, but scaling remains. Training-window
means may weight overlapping raw samples repeatedly; they need not equal
means computed on unique raw samples. The fit population must be recorded.
\paragraph{Why ECDR needs it.}
The reference is the composition of preprocessing and CNN, not the CNN
weights alone. A certificate in normalized coordinates must be checked in
those coordinates or transported correctly to raw space. Replacing fitted
scales changes the function being certified.
\paragraph{Conditions and connections.}
Positive finite scales, finite inputs and a frozen fit are assumptions.
Per-window centering defines a different map and can destroy exact
normalized overlap. This operation does not remove gravity or independently
calibrate a sensor. B08 derives the corresponding transport of bounds.

\concept{B08}{Transporting a metric, a region and an affine bound}{F16, F20, B07}
\paragraph{Foundation and definition.}
Coordinate transport applies the same invertible map to every object used
in a claim. With $z=D^{-1}(x-\mu)$ and fixed positive diagonal $D$,
\[
 \norm{z-z^r}_2^2=\sum_{i=1}^{d}\frac{(x_i-x_i^r)^2}{\sigma_i^2}.
\]
\paragraph{Symbols and meaning.}
$d$ is input dimension; $\mu\in\R^d$ is the fixed fitted mean;
$x,x^r\in\R^d$ are raw query and anchor;
$z,z^r$ are their normalized images; $\sigma_i=D_{ii}>0$ is a scale.
The normalized norm is dimensionless if scales carry the corresponding
raw units. Raw channels in g and rad/s have no shared physical unit merely
because their numeric values are added in a norm.
\paragraph{Derivation and example.}
A raw box with budgets $b_i^{\rm raw}\ge0$ and $|x_i-x_i^r|\le b_i^{\rm raw}$ maps to
$|z_i-z_i^r|\le b_i^{\rm raw}/\sigma_i$. If normalized margin
$\gamma$ has a valid bound $\gamma(z)\ge a_z^\top z+\beta_z$, substitution gives
\[
 a_x=D^{-\top}a_z,\qquad
 \beta_x=\beta_z-a_z^\top D^{-1}\mu,
 \qquad g(x)=\gamma(\mathcal P(x))\ge a_x^\top x+\beta_x.
\]
$a_z,a_x\in\R^d$ are bound coefficients and $\beta_z,\beta_x\in\R$ are
intercepts. For scalars $a_z=3,\beta_z=1,\mu=2,\sigma=0.5$, obtain
$a_x=6,\beta_x=-11$; both expressions equal 1.6 at raw $x=2.1$.
\paragraph{Why ECDR needs it.}
Using a normalized radius against raw distances accepts the wrong set.
Transport is the link between a mathematical domain and the physical
observations admitted by a checker. Changing units can change a raw
numeric distance while representing the same physical variation.
\paragraph{Conditions and connections.}
The lower bound must hold on the image of the raw domain. Invertibility
alone does not make a bound globally valid. C14 applies the same reasoning
to projected and generator-based regions.

\concept{B09}{The shift--append operation is a linear state update}{F15, F17, B03}
\paragraph{Foundation and definition.}
A window advances by dropping its first $s$ columns and appending a new
wave. Define
\[
 \mathcal T_s(X,U)=S_sX+J_sU=[X[:,s:W],U].
\]
\paragraph{Symbols and meaning.}
$X\in\R^{C\times W}$ is the old window, $U\in\R^{C\times s}$ is new
data, and $1\le s\le W$. $S_s$ means drop-and-zero-fill:
$S_sX=[X[:,s:W],0_{C\times s}]$.
$J_sU=[0_{C\times(W-s)},U]$ inserts the new tail. They are linear maps
between the stated tensor spaces; $0_{a\times b}$ is an all-zero matrix.
The brackets concatenate columns, not add them.
\paragraph{Derivation and example.}
With $C=1,W=4,s=2$, vectorized matrices are
\[
 S_s=\begin{pmatrix}0&0&1&0\\0&0&0&1\\0&0&0&0\\0&0&0&0\end{pmatrix},
 \qquad J_s=\begin{pmatrix}0&0\\0&0\\1&0\\0&1\end{pmatrix}.
\]
For $X=(a,b,c,d)^\top$ and $U=(u,v)^\top$, their sum is
$(c,d,u,v)^\top$. $J_s$ is injective: reading its last $s$ positions
recovers $U$, so two different tails cannot create the same full update
when the old window is fixed.
\paragraph{Why ECDR needs it.}
This operation exposes exactly which variables are fixed and which are
new. A one-step wave region can vary $Cs$ coordinates instead of all
$CW$ coordinates. It also defines a state transition for multi-step proofs.
\paragraph{Conditions and connections.}
The update assumes the same recording, channel order and sample grid.
Dropped samples are lost from this state; $S_s$ is generally not invertible.
Reducing free input dimension does not by itself prove lower verifier cost.
B10 identifies the equality inherited by shared samples.

\concept{B10}{Aligned overlap is an identity of sample indices}{F09, B04, B09}
\paragraph{Foundation and definition.}
Two consecutive windows from the same fixed grid share samples. For
recording $v[n]\in\R^C$ and $X_k[:,i]=v[ks+i]$,
\[
 X_{k+1}[:,i]=X_k[:,i+s],\qquad 0\le i<W-s.
\]
\paragraph{Symbols and meaning.}
$k$ is window index, $i$ is a within-window sample index, $s$ is stride,
$W$ is window length and $C$ is channel count, with $1\le s\le W$ here. The same vector appears at
different tensor positions: new position $i$ and old position $i+s$.
There are $C(W-s)$ shared scalar coordinates after alignment.
\paragraph{Derivation and example.}
Substitute the window definition:
$X_{k+1}[:,i]=v[(k+1)s+i]=v[ks+(i+s)]$.
The final expression is the definition of $X_k[:,i+s]$. This proof uses
index equality, not smoothness, a common activity label or high correlation.
For $C=6,W=128,s=64$, 384 shared values remain and 384 new values arrive.
Time-homogeneous per-channel normalization preserves equality: the same
channel map is applied at every temporal position. Separately fitted window means need
not preserve it.
\paragraph{Why ECDR needs it.}
Exact feature sharing for a valid CNN starts from this identity. It can
also justify checking only a new tail if the old overlap is guaranteed by
the acquisition contract. A checker without that guarantee must verify the
overlap too. Data similarity alone does not supply sample identity.
\paragraph{Conditions and connections.}
There is no such guarantee across recording boundaries, dropped samples,
changed channel layouts or asynchronous resampling. B11 shows why this
identity does not imply a small ordinary difference; B24 lifts it through
appropriate convolutional layers.

\concept{B11}{Same-coordinate perturbation is not aligned overlap}{F20, B09--B10}
\paragraph{Foundation and definition.}
An ordinary difference compares matching coordinate labels in two input
tensors, not the matching raw samples. For old $X$ and new
$X'=\mathcal T_s(X,U)$,
\[
 \Delta=X'-X=(S_s-I)X+J_sU.
\]
\paragraph{Symbols and meaning.}
$X,X',\Delta\in\R^{C\times W}$; $U\in\R^{C\times s}$ is new data;
$S_s,J_s$ are B09's linear maps; $I$ is the identity on full windows.
The Frobenius norm $\norm{\Delta}_F$ is Euclidean norm after vectorization.
It aggregates differences at the same channel and window position.
\paragraph{Derivation and counterexample.}
For one channel, $W=4,s=2$, take $X=(0,0,M,M)$ and $U=(0,0)$, where
$M\in\R$ is arbitrary amplitude. Then $X'=(M,M,0,0)$ still satisfies exact
50\% aligned overlap, but
\[
 \norm{X'-X}_2=\sqrt{M^2+M^2+M^2+M^2}=2|M|.
\]
No small upper bound independent of amplitude follows from overlap. If an
additional assumption imposes $|X[c,i]|,|U[c,j]|\le A$ for finite $A\ge0$,
each difference is at most $2A$, giving
$\norm{\Delta}_F\le2A\sqrt{CW}$. This bound comes from amplitude control.
\paragraph{Why ECDR needs it.}
An ordinary norm ball may need to include changes due to moving old samples
as well as new data. Reusing the last class merely because half the samples
are shared is unjustified. Wave geometry can encode this structure directly.
\paragraph{Conditions and connections.}
The example disproves a logical implication, not an empirical claim about
human movement. Bounds on input distance still require margin sensitivity
to obtain a decision certificate, as in C02.

\concept{B12}{Marginal coverage, joint coverage and the union bound}{F29--F34, F24}
\paragraph{Foundation and definition.}
A box accepts a query only when every coordinate constraint holds. Let
$\Delta=(\Delta_1,\ldots,\Delta_d)$ be a random perturbation and define
$E_i=\{|\Delta_i|\le b_i\}$ for fixed budgets $b_i\ge0$.
Box membership is the joint event $\bigcap_{i=1}^dE_i$.
\paragraph{Symbols and meaning.}
$d\ge1$ is dimension, $i$ indexes coordinates, $\Pr$ is probability under
a specified data-generating distribution, and $\alpha_i\in[0,1]$ is a
valid upper bound on $\Pr(E_i^c)$. The complement $E_i^c$ is coordinate
failure. A train-set quantile does not automatically provide this
population probability guarantee.
\paragraph{Derivation and example.}
Countable subadditivity from the probability axioms gives
\[
 \Pr\!\left(\bigcap_iE_i\right)
 =1-\Pr\!\left(\bigcup_iE_i^c\right)
 \ge\max\left(0,1-\sum_i\alpha_i\right).
\]
No independence is used. If, additionally, the events are independent and
each has probability $q$, their joint probability is $q^d$.
For $q=0.9,d=768$, it is approximately $7.21518\times10^{-36}$ under
that added assumption. This is a constructed example, not a HAPT estimate.
With two identical events of probability 0.9, joint coverage is instead 0.9.
\paragraph{Why ECDR needs it.}
Per-coordinate q90 budgets do not establish q90 joint membership, much
less q90 certified reuse. Geometry, empirical joint coverage and a
whole-domain model property are three separate questions.
\paragraph{Conditions and connections.}
The distribution and fit/evaluation split must be declared before making
statistical claims. A certificate can be distribution-free conditional on
membership; its usefulness still depends on how often the stream enters
the domain. C18 and C28 separate those claims.

\concept{B13}{Temporal dependence and the unit of a split}{F29--F33, B10}
\paragraph{Foundation and definition.}
Independence means a joint probability factors appropriately; it is not a
consequence of storing observations in different tensor rows. Overlapping
windows contain common raw samples. Multiple recordings from one person
can also share person-specific behavior and device placement.
\paragraph{Symbols and meaning.}
Let $Y_1,\ldots,Y_N$ be scalar evaluation outcomes with finite second
moments, $N\ge1$ their count, and $\bar Y=N^{-1}\sum_iY_i$ their mean.
$\operatorname{Cov}(Y_i,Y_j)$ is covariance. Then
\[
 \operatorname{Var}(\bar Y)=\frac1{N^2}
 \sum_{i=1}^{N}\sum_{j=1}^{N}\operatorname{Cov}(Y_i,Y_j).
\]
\paragraph{Derivation and example.}
Expand the centered square of the sum and use expectation linearity.
The familiar $\sigma^2/N$ requires zero off-diagonal covariances and
appropriate common variance $\sigma^2$. If $Y_i=Y_1$ for every $i$,
averaging duplicate rows gives variance $\operatorname{Var}(Y_1)$, not
$\operatorname{Var}(Y_1)/N$. This extreme example explains why row count
alone does not measure independent information.
\paragraph{Why ECDR needs it.}
A random window split may put the same raw sample into training and test.
Splitting by subject avoids that particular overlap but does not make all
within-subject queries independent. Reuse should be replayed in recording
order because its anchor state depends on preceding queries.
\paragraph{Conditions and connections.}
The statistical unit, clustering and dependence assumptions must match the
uncertainty claim. A subject-level bootstrap is a possible design, not an
already-performed experiment here. Describing all splits for exploration
must not be called an untouched confirmatory test afterward.

\concept{B14}{The reference classifier is a composition, not just weights}{F04, F17, B07}
\paragraph{Foundation and definition.}
The reference function maps a raw query into class scores through a fixed
preprocessing map followed by a fixed network:
\[
 F=f_\theta\circ\mathcal P:\R^d\to\R^K,
 \qquad F(x)=f_\theta(\mathcal P(x)).
\]
\paragraph{Symbols and meaning.}
$d=CW$ is raw input dimension, $K\ge2$ is class count,
$x\in\R^d$ is one vectorized raw window, $\mathcal P$ is preprocessing,
and $f_\theta$ is the CNN with parameters and frozen operator state
denoted collectively by $\theta$. Its input is normalized vector $z$;
$F$'s input is raw vector $x$. Class ordering and tie handling belong to
the decision contract even though neither appears as a trainable weight.
\paragraph{Derivation and example.}
Let scalar preprocessing be $\mathcal P(x)=2x-4$ and an illustrative
two-score network be $f_\theta(z)=(z,-z)$. Then
$F(x)=(2x-4,4-2x)$. The winner changes at raw value $x=2$, while the
network boundary is normalized value $z=0$. Substituting raw $x$ directly
into the network uses a different classifier.
\paragraph{Why ECDR needs it.}
A reusable certificate must match the graph, weights, preprocessing,
class schema, deterministic mode and domain coordinates used at runtime.
A matching filename is not evidence of matching function identity.
Certification relative to a reference means preserving its decision,
including occasions when that reference disagrees with ground truth.
\paragraph{Conditions and connections.}
Real-valued functional identity and floating-point program identity are
distinct obligations. Updating weights or fitted normalization requires
new evidence unless an explicit transfer result applies. C20 handles
numerical semantics; C23 handles anchor identity through time.

\concept{B15}{Logits, maximum, argmax and deterministic ties}{F06, F45, B14}
\paragraph{Foundation and definition.}
A logit vector is a list of real scores. For $F(x)\in\R^K$, the classifier
returns an index with largest score:
\[
 h(x)=\operatorname{argmax}_{1\le j\le K}F_j(x).
\]
\paragraph{Symbols and meaning.}
$x\in\R^d$ is a query, $F_j(x)\in\R$ is the score for class index $j$,
$K\ge2$ is class count and $h(x)\in\{1,\ldots,K\}$ is the selected
class. A deterministic tie rule makes $h$ a function rather than a
set-valued relation. In contrast, $\max_jF_j(x)$ returns a score value,
not the class index. Logits do not have sensor units or automatic
probabilistic calibration.
\paragraph{Derivation and example.}
For scores $(2.4,1.1,2.0)$, the maximum is 2.4 and the unique argmax is
class~1. For $(2,2,0)$, the set of maximizers is $\{1,2\}$; a
smallest-index rule returns1. Softmax at temperature $\tau>0$ gives
$p_j=\exp(F_j/\tau)/\sum_k\exp(F_k/\tau)$, where $k$ is the denominator
index. Its positive common denominator and the strict monotonicity of
exponentiation preserve score ordering in exact real arithmetic.
\paragraph{Why ECDR needs it.}
Class-only reuse need not recreate current logits or probabilities. A
downstream consumer requiring an embedding or calibrated confidence asks
for a different output contract. Margin certification can therefore work
directly on logits without inserting softmax.
\paragraph{Conditions and connections.}
Strict positive pairwise margins avoid relying on ties, but are sufficient
rather than necessary for a particular deterministic tie rule. NaN scores
are outside the finite-real model. B16 formalizes the pairwise test;
C20 accounts for numerical errors in actual programs.

\concept{B16}{A pairwise decision margin converts argmax into inequalities}{F06, B15}
\paragraph{Foundation and definition.}
For a fixed target class $c$, compare its score with each competing class.
The pairwise margin is
\[
 g_{cj}(x)=F_c(x)-F_j(x),\qquad j\in\{1,\ldots,K\}\setminus\{c\}.
\]
\paragraph{Symbols and meaning.}
$x\in\R^d$ is raw input; $F:\R^d\to\R^K$ is the frozen reference;
$c$ is a chosen target index and $j$ is a different competitor. Each
$g_{cj}(x)\in\R$ is a scalar score difference. For $K=12$, there are
11 separate comparisons. Define the smallest target margin at an anchor
$x^r$ as $m_r=\min_{j\ne c}g_{cj}(x^r)$.
\paragraph{Derivation and example.}
If every $g_{cj}(x)>0$, then $F_c(x)>F_j(x)$ for every competitor, so
$h(x)=c$ uniquely. For $(F_1,F_2,F_3)=(4.5,-0.5,1)$ and $c=1$,
the two margins are 5 and 3.5. At another input with scores
$(6.5,-2.5,5)$, target score increases but the smaller margin falls to
1.5. Target confidence cannot be inferred from its own logit alone.
\paragraph{Why ECDR needs it.}
Reuse requires preserving ordering, not preserving every score numerically.
Bounding a difference directly can retain cancellations between two logits
that are lost when each logit is bounded independently. A verifier must
check all competitors, not just the class that is second-largest at the
anchor: a different competitor can become dangerous elsewhere in a region.
\paragraph{Conditions and connections.}
Strict positivity is sufficient; equality depends on the tie rule.
A positive margin at one input is a pointwise fact, not a whole-domain
certificate. B17 supplies a local existence theorem from continuity;
C18 states the universal lower-bound condition used by ECDR.

\concept{B17}{Local invariance from continuity and a positive anchor margin}{F26, B16}
\paragraph{Foundation and definition.}
A classifier can be locally constant in its decision even though its
logits vary continuously. Suppose $F:\R^d\to\R^K$ is continuous at
anchor $x^r$, $K$ is finite, and the chosen class $c$ has
$g_{cj}(x^r)>0$ for every $j\ne c$.
\paragraph{Symbols and meaning.}
$g_{cj}=F_c-F_j$ is a pairwise margin; $d$ is input dimension;
$\norm{\cdot}$ is a fixed norm; and $\epsilon_j>0$ will denote an
input radius for competitor $j$. A ball of radius $\epsilon$ contains
all $x$ satisfying $\norm{x-x^r}<\epsilon$.
\paragraph{Derivation.}
Each margin is continuous because it is a difference of continuous
functions. For each $j$, continuity provides $\epsilon_j>0$ such that
\[
 \norm{x-x^r}<\epsilon_j
 \ \Longrightarrow\
 |g_{cj}(x)-g_{cj}(x^r)|<g_{cj}(x^r)/2.
\]
Consequently $g_{cj}(x)>g_{cj}(x^r)/2>0$. Take
$\epsilon=\min_{j\ne c}\epsilon_j$. A minimum of finitely many positive
numbers is positive. Thus every input in that ball has unique winner $c$.
\paragraph{Why ECDR needs it.}
The theorem explains why a certified invariant neighborhood is possible
around a strictly classified input. It does not compute its useful size.
Operational bounds such as Lipschitz or affine relaxations turn this
existence statement into checkable sufficient conditions.
\paragraph{Conditions and connections.}
Finite $K$, continuity at the anchor and strict margins are necessary
assumptions for this proof. It says neither that the radius is independent
of the model nor that a future query enters the ball. B18 shows the former
failure; C02 constructs a quantitative radius under sensitivity bounds.

\concept{B18}{Why no positive radius can be universal over all classifiers}{F02, B16--B17}
\paragraph{Foundation and definition.}
To disprove a universal claim, construct one permitted counterexample for
each proposed bound. Consider the claim that some fixed radius $r_0>0$
preserves the decision of every continuous classifier at every strictly
classified anchor. The claim reverses the order of quantifiers in B17.
\paragraph{Symbols and meaning.}
$r_0\in\R$ is a proposed positive universal radius. Define
$\varepsilon=r_0/2$, scalar input $x\in\R$, anchor $x^r=0$, and a
two-class reference
\[
 F^{(\varepsilon)}(x)=(\varepsilon-x,0).
\]
The superscript labels the chosen model, not exponentiation. Its target
class is $c=1$ at the anchor.
\paragraph{Derivation and counterexample.}
The margin is $g_{12}(x)=\varepsilon-x$, a continuous affine function
with Lipschitz constant~1. At the anchor its margin is $\varepsilon>0$.
Take $x'=3r_0/4$. Then $|x'-x^r|=3r_0/4<r_0$, but
\[
 g_{12}(x')=r_0/2-3r_0/4=-r_0/4<0.
\]
Class2 wins within the proposed radius. This construction works for every
$r_0>0$. It does not contradict a model-dependent radius smaller than
$\varepsilon$.
\paragraph{Why ECDR needs it.}
Input overlap, network continuity or a domain's descriptive name cannot
justify a fixed reusable radius for arbitrary weights and margins.
The region must be linked to a particular reference, anchor and bound.
\paragraph{Conditions and connections.}
This is a distribution-free logical counterexample, not a claim that HAPT
models behave this way. Even a known Lipschitz constant is insufficient
without an anchor margin. C02 makes that dependence explicit.

\concept{B19}{Convolution as a sparse affine map with shared coefficients}{F09, F15--F16, B03}
\paragraph{Foundation and definition.}
The Conv1D convention used here is cross-correlation: kernel coefficients
are not reversed. For valid padding and stride1,
\[
 V[o,i]=b_o+\sum_{c=0}^{C_{\rm in}-1}\sum_{r=0}^{k-1}
 W[o,c,r]Z[c,i+r].
\]
\paragraph{Symbols and meaning.}
$Z\in\R^{C_{\rm in}\times T}$ is input, $V\in\R^{C_{\rm out}\times(T-k+1)}$
is preactivation, $W\in\R^{C_{\rm out}\times C_{\rm in}\times k}$ is
the kernel and $b\in\R^{C_{\rm out}}$ is bias. $o,c,r,i$ index output
channel, input channel, kernel tap and output position. $1\le k\le T$.
The same $W[o,c,r]$ is used at all output positions.
\paragraph{Derivation and example.}
For one channel with $Z=(1,2,3,4)$, kernel $(1,-1)$ and bias~0.5,
each output is $Z[i]-Z[i+1]+0.5=-0.5$, giving three outputs. Stacking
all output coordinates gives $\operatorname{vec}(V)=\mathsf W\operatorname{vec}(Z)+\mathsf b$.
Writing $T_o=T-k+1$, matrix $\mathsf W$ has shape
$(C_{\rm out}T_o)\times(C_{\rm in}T)$, and $\mathsf b\in\R^{C_{\rm out}T_o}$
repeats channel biases. The sparse matrix's repeated entries encode weight
sharing. This is affine, not generally linear when bias is nonzero.
\paragraph{Why ECDR needs it.}
The affine representation allows interval sign rules and linear
relaxation methods to treat convolution rigorously. Local shared weights
also support exact shifted feature equality, but only at correctly aligned
positions whose receptive fields contain equal inputs.
\paragraph{Conditions and connections.}
Stride, dilation, grouping and padding change connectivity and output
length; they must be stated rather than inferred from the word CNN.
The main theoretical model uses valid stride~1/dilation~1 convolutions.
The separate English CNN document gives an individual page to each
operator parameter. B24 proves the corresponding feature identity.

\concept{B20}{ReLU creates affine pieces without preserving every margin}{F16--F17, B19}
\paragraph{Foundation and definition.}
The rectified linear unit is the pointwise scalar function
$\relu(v)=\max(0,v)$. Applied to tensor $V$, it gives $H[o,i]=\relu(V[o,i])$
without changing shape or introducing trainable parameters.
\paragraph{Symbols and meaning.}
$v\in\R$ is one preactivation, $V$ is a preactivation tensor, $H$ is
its activation tensor, and $o,i$ index a feature channel and position.
For $m\ge1$ neurons, let $w\in\R^m$ be their preactivation vector.
On a fixed activation pattern, define $D\in\R^{m\times m}$ diagonal
with entries~0 or~1. Then $\relu(w)=Dw$. Neurons identically zero may use a zero gate.
\paragraph{Derivation and example.}
If every preactivation sign is fixed on an input region, replacing each
ReLU by its fixed gate makes the full affine/ReLU network affine there:
composition of affine maps is affine. For a scalar affine preactivation
$v(x)=2x-1$, ReLU is~0 when $x\le1/2$ and $2x-1$ when $x\ge1/2$.
The pieces agree at the boundary but have different slopes.
For any scalar $a,b$, case analysis gives
$|\relu(a)-\relu(b)|\le|a-b|$; coordinatewise aggregation yields the
same nonexpansion for $L_1,L_2,L_\infty$ norms.
\paragraph{Why ECDR needs it.}
Stable activation patterns allow exact affine reasoning; unstable neurons
require sound bounds. A pattern observed only at the anchor is not
established over an entire candidate region. Network layers around ReLU
may amplify differences even though ReLU alone is nonexpansive.
\paragraph{Conditions and connections.}
Nonnegative activations do not imply increasing class margins: head
coefficients can be negative. C04 propagates intervals; C05 derives valid
affine bounds for a neuron whose sign can change.

\concept{B21}{Receptive-field span and jump on the input clock}{B01--B02, B19}
\paragraph{Foundation and definition.}
An architectural receptive field describes input positions that can affect
one feature. Its span is the distance between extreme potentially accessed
indices plus1. The jump is the input-grid displacement between consecutive
feature positions. This temporal extent is not a vector-space linear span.
\paragraph{Symbols and meaning.}
For local layer $\ell$, let $k_\ell\ge1$ be kernel taps, $d_\ell\ge1$
dilation and $a_\ell\ge1$ stride. Let $R_\ell$ be receptive-field span
in samples and $J_\ell$ input-grid jump in samples. With
$R_0=J_0=1$,
\[
 R_\ell=R_{\ell-1}+(k_\ell-1)d_\ell J_{\ell-1},\qquad
 J_\ell=a_\ell J_{\ell-1}.
\]
\paragraph{Derivation and example.}
The extreme input features in the next kernel are
$(k_\ell-1)d_\ell$ positions apart. Multiply by their raw-grid spacing
$J_{\ell-1}$ and add one preceding field's extent. Advancing the output
by one moves $a_\ell$ preceding positions, giving the jump recurrence.
With kernels~5,3,3 and all strides/dilations~1, spans are~5,7,9 and jumps
remain1. At 50 Hz, RF~9 has first-to-last interval8/50=0.16 s, while
sample-count duration is9/50=0.18 s.
\paragraph{Why ECDR needs it.}
Feature equality requires the whole relevant receptive field to lie in
the shared portion. Input overlap~50\% therefore does not mean feature
overlap~50\%. Padding and striding require extra boundary and phase reasoning.
\paragraph{Conditions and connections.}
With dilation, the span may contain holes; zero learned weights may reduce
actual dependence. The recurrence describes potential connectivity, not
empirical feature importance. B24 specializes it to a contiguous valid CNN.

\concept{B22}{Global average pooling and the affine classifier head}{F09, F16, B20}
\paragraph{Foundation and definition.}
Global average pooling (GAP) reduces a final feature map along time while
retaining each channel. For normalized input $Z$, define
\[
 m(Z)=\frac1n\sum_{i=0}^{n-1}H_i(Z),\qquad
 f_\theta(Z)=A\,m(Z)+\beta.
\]
\paragraph{Symbols and meaning.}
$H_i(Z)\in\R^{d_h}$ is the vector of final features at position $i$,
$n\ge1$ is the number of final positions, $d_h$ is feature-channel count,
$m(Z)\in\R^{d_h}$ is the pooled vector, $A\in\R^{K\times d_h}$ is
the head matrix and $\beta\in\R^K$ is its bias. Raw logits are
$F(x)=f_\theta(\mathcal P(x))$. Rows $A_{j,:}$ correspond to classes.
\paragraph{Derivation and example.}
For feature positions $(2,1),(3,1),(4,1),(5,1)$,
$m=(14/4,4/4)=(3.5,1)$. Averaging is linear in the feature map, so
equal channel sums give equal pooled vectors. It discards the explicit
ordering of feature positions at the head. It does not make the full CNN
invariant to arbitrary permutations of raw samples: the local features
themselves can change under such permutations.
\paragraph{Why ECDR needs it.}
Feature sums allow shared contributions to cancel exactly when a window
shifts. An affine head then turns a pooled change into a direct pairwise
margin change. A position-weighted Flatten--MLP head does not generally
have the same cancellation.
\paragraph{Conditions and connections.}
Equal position counts and the stated pooling/head operators are essential.
A nonlinearity after pooling needs separate analysis. B25 derives the GAP
update, and B26 connects it to decision certification rather than just
activation reuse.

\concept{B23}{Parameter count and arithmetic work are different quantities}{B19, B22}
\paragraph{Foundation and definition.}
Parameters count stored learnable scalars. Arithmetic work counts uses of
those scalars while evaluating one input. For a bias-enabled, groups~1
convolution with input channels $C_{\rm in}$, output channels $C_{\rm out}$,
kernel length $k$ and output length $T_{\rm out}$,
\[
 P=C_{\rm out}(C_{\rm in}k+1),\qquad
 M=T_{\rm out}C_{\rm out}C_{\rm in}k.
\]
\paragraph{Symbols and meaning.}
$P$ is parameter count, $M$ is multiply--accumulate (MAC) count under one
product-and-accumulation convention. The added1 in $P$ is one bias per
output channel; it is not another kernel tap. Output length multiplies
work because weights are reused across positions, not because new weights
are created. A $d_h$-to-$K$ affine head has $K(d_h+1)$ parameters and
$Kd_h$ MAC under the same convention.
\paragraph{Derivation and example.}
For valid kernels~5,3,3, channels $6\to32\to64\to64$ and input length~128,
output lengths are124,122,120. The three parameter counts are~992,6208,
12352; a 12-class head adds~780, totaling~20332. MAC counts are~119040,
749568,1474560 and768, totaling~2343936. Reporting two FLOP/MAC gives
4687872 FLOP for these operators only.
\paragraph{Why ECDR needs it.}
MAC reduction suggests computational opportunity but is not a latency
theorem. ReLU, GAP, bias additions, normalization, memory traffic, dispatch
and certificate work still cost time. C26--C27 account for that overhead.
\paragraph{Conditions and connections.}
Precision, kernel algorithms, buffer schedules and hardware affect actual
runtime and peak memory. Float32 weight payload~20332$\times$4=81328 bytes
is not checkpoint file size or process RAM. These are theoretical counts,
not measured quality or speed of a new HAPT model.

\concept{B24}{Exact aligned feature equality through a valid CNN}{B10, B19--B21}
\paragraph{Foundation and definition.}
Shared input samples induce shared features only under an operator contract.
Assume frozen time-homogeneous per-channel preprocessing, valid stride~1/dilation~1 convolutions,
position-independent pointwise activations, fixed weights and no pooling
or hidden state. Let the final receptive-field span be $R$ and final
length be $n=W-R+1$.
\paragraph{Symbols and meaning.}
$Z,Z'\in\R^{C\times W}$ are consecutive normalized windows shifted by
$s$ samples, with $1\le s\le n$;
$Z'[:,i]=Z[:,i+s]$ for $0\le i<W-s$.
$H_i(Z)\in\R^{d_h}$ is the final feature at position $i$.
The claim is
\[
 H_i(Z')=H_{i+s}(Z),\qquad 0\le i<n-s.
\]
\paragraph{Derivation.}
The new feature at $i$ reads input indices $i,\ldots,i+R-1$.
For $i\le n-s-1$, the last index is at most $W-s-1$, inside overlap.
Each sample there equals the old sample at offset $s$. Apply equal
kernels, biases and pointwise activations layer by layer to obtain equal
outputs. The proof establishes shifted equality, not equality at the same
tensor position.
\paragraph{Why ECDR needs it.}
For $W=128,R=9,s=64$, final length~120 leaves~56 shared positions.
The~64 new final features depend on~72 raw columns:~64 new samples plus
an~8-sample halo. This explains why forwarding only the~64 newly arrived
raw columns through the entire valid backbone misses needed features.
\paragraph{Conditions and connections.}
Padding, striding, per-window preprocessing or position-dependent operators
require a revised proof. Equality of features alone does not fix the
new pooled sum or class. B25--B26 derive the remaining changes.

\concept{B25}{The GAP update is a cancellation identity}{F09, B22, B24}
\paragraph{Foundation and definition.}
When final feature maps share a shifted interior, their sums differ only
in dropped and added positions. Define
\[
 D_{\rm drop}=\sum_{i=0}^{s-1}H_i(Z),\qquad
 D_{\rm new}(U)=\sum_{i=n-s}^{n-1}H_i(Z'),
 \qquad Z'=\mathcal T_s(Z,U).
\]
\paragraph{Symbols and meaning.}
$Z,Z'\in\R^{C\times W}$ are old/new normalized windows; $U\in\R^{C\times s}$
is the new wave; $n$ is final feature length; $1\le s\le n$ is shift;
$H_i\in\R^{d_h}$ is a feature vector. Both $D_{\rm drop}$ and
$D_{\rm new}$ are $d_h$-dimensional sums. GAP is
$m(Z)=n^{-1}\sum_iH_i(Z)$.
\paragraph{Derivation and example.}
By B24, the shared new sum equals $\sum_{i=s}^{n-1}H_i(Z)$. Therefore
\[
 \sum_iH_i(Z')-\sum_iH_i(Z)=D_{\rm new}(U)-D_{\rm drop},
 \quad m(Z')-m(Z)=\frac{D_{\rm new}(U)-D_{\rm drop}}n.
\]
For old features $(0,1),(1,1),(2,1),(3,1)$ and new features
$(2,1),(3,1),(4,1),(5,1)$, take $n=4,s=2$.
The two sums are $(1,2)$ and $(9,2)$, giving pooled delta $(2,0)$.
\paragraph{Why ECDR needs it.}
Cancellation isolates the unknown future contribution. If that contribution
can be bounded over a wave region, the classifier may be certified before
computing exact features for an observed query. Computing its actual sum
instead gives an incremental evaluation, a different form of reuse.
\paragraph{Conditions and connections.}
This is exact real algebra under B24's assumptions and equal feature
lengths. Floating-point sum updates may not match recomputation bitwise.
An identity provides no bound on all possible new waves by itself;
B26 and C19 provide the margin connection.

\concept{B26}{A pooled-feature change induces a pairwise-margin change}{B16, B22, B25}
\paragraph{Foundation and definition.}
An affine head turns a GAP change into an exact scalar margin change.
Define the normalized-network margin
\[
 \gamma_{cj}(Z)=f_{\theta,c}(Z)-f_{\theta,j}(Z),\qquad
 q_j=(A_{c,:}-A_{j,:})^\top.
\]
\paragraph{Symbols and meaning.}
$A\in\R^{K\times d_h}$ and $\beta\in\R^K$ are head weights and bias;
$c$ is target class, $j\ne c$ is a competitor, and $q_j\in\R^{d_h}$
is a column vector. $Z$ and $Z'$ are old/new normalized inputs;
$n$ is feature length. For raw query $x$, $g_{cj}(x)=\gamma_{cj}(\mathcal P(x))$.
The new wave is $U\in\R^{C\times s}$, where $C$ is input-channel count and
$1\le s\le n$ is the shift. Vectors $D_{\rm new}(U),D_{\rm drop}\in\R^{d_h}$
sum newly added and dropped feature positions, respectively (B25).
\paragraph{Derivation and example.}
Using the head formula and subtracting old from new cancels its bias:
\[
 \gamma_{cj}(Z')=\gamma_{cj}(Z)+
 \frac1nq_j^\top\big(D_{\rm new}(U)-D_{\rm drop}\big).
\]
The bias still contributes to the old margin; it must not be removed
from that value. In B25's example, pooled delta is $(2,0)$.
For $q_2=(2,-1)^\top$ the margin increases~4; for
$q_3=(-1,1)^\top$ it decreases~2. The same feature change helps one
comparison and harms another.
\paragraph{Why ECDR needs it.}
This formula expresses exactly what needs a bound for class preservation.
Bounding every feature or logit independently can lose cancellation;
bounding the pairwise scalar objective can be more informative. The
whole-domain lower bound must remain valid for all admitted waves.
\paragraph{Conditions and connections.}
The identity needs the specified affine head and B25's feature update.
It is neither an empirical stability estimate nor a certificate until
the unknown term is enclosed. C19 chooses the correct feature-sum bounds
according to the signs of the coefficients in $q_j$.

\concept{B27}{A complete miniature CNN: forward arithmetic}{B19--B22, B26}
\paragraph{Foundation and definition.}
A small fully specified network makes each stage independently checkable.
Take identity preprocessing, one channel, window length~6, shift~2, a
valid length~3 convolution, two filters, ReLU, GAP and a three-class head.
\paragraph{Symbols and meaning.}
The old input is $X=(1,2,3,4,5,6)$, the appended wave is $U=(7,8)$,
and $X'=(3,4,5,6,7,8)$. Filter weights and biases are
\[
 w_1=(1,-1,1),\ b_1=-2;\qquad w_2=(-1,0,1),\ b_2=-1.
\]
The head matrix $A\in\R^{3\times2}$ and bias $\beta\in\R^3$ are
\[
 A=\begin{pmatrix}1&-1\\-1&0\\2&-2\end{pmatrix},\qquad
 \beta=(4,1,0)^\top.
\]
\paragraph{Derivation.}
Filter~1 on the first old patch gives $1-2+3-2=0$; filter~2 gives
$-1+3-1=1$. The position-wise feature vectors after ReLU are
\[
 H(X)=((0,1),(1,1),(2,1),(3,1)),\quad
 H(X')=((2,1),(3,1),(4,1),(5,1)).
\]
Thus $m(X)=(3/2,1)$, $m(X')=(7/2,1)$ and the raw/normalized logits,
identical here because preprocessing is identity, are
\[
 F(X)=(9/2,-1/2,1),\qquad F(X')=(13/2,-5/2,5).
\]
Both choose class~1; its margins change from $(5,7/2)$ to $(9,3/2)$.
\paragraph{Why ECDR needs it.}
The example verifies score ordering, feature alignment, GAP cancellation
and opposite margin responses in a model whose every operation is visible.
It illustrates a mechanism rather than demonstrating HAPT accuracy.
\paragraph{Conditions and connections.}
All values here are exact rational arithmetic. No training or data fit is
claimed. One observed wave is not a region certificate; B28 varies it over
a continuous box and proves the necessary signs.

\concept{B28}{Certifying a continuous wave box in the miniature CNN}{B27, C03, C18--C19}
\paragraph{Foundation and definition.}
Replace B27's single wave by $U=(u,v)$ in the continuous region
$\mathcal U=[6.5,7.5]\times[7.5,8.5]$. The old anchor and network remain
fixed. Every new input is $X'=(3,4,5,6,u,v)$.
\paragraph{Symbols and meaning.}
$u,v\in\R$ are the two new samples; $\mathcal U$ is a set, not an
assumed probability distribution. There are $n=4$ final positions;
$D_{\rm new}\in\R^2$ sums the two new feature vectors.
The two class~1 margins are $g_{12}$ and $g_{13}$, with competitors~2 and~3.
\paragraph{Derivation.}
The new tail preactivations are
\[
 V_2=(u-3,u-6),\qquad V_3=(4-u+v,v-7).
\]
Their coordinate lower bounds are $(3.5,0.5)$ and $(4,0.5)$, all positive.
Thus post-ReLU features satisfy $H_2=V_2,H_3=V_3$ throughout the box,
not just at its center.
Summing symbolically gives
\[
 D_1=1+v,\quad D_2=u+v-13,\quad
 g_{12}=5+\frac{15+v-u}{4},\quad g_{13}=\frac{u-1}{4}.
\]
$D_1,D_2$ are feature-sum coordinates. Optimizing these affine margins
over the stated intervals yields $\min g_{12}=35/4$ and
$\min g_{13}=11/8$, both positive. Every wave in the box preserves class~1.
\paragraph{Why ECDR needs it.}
This is an actual whole-domain argument without assuming independent or
uniform samples. Keeping shared variables in the objective retains
dependencies. Bounding tail positions independently gives the looser but
still sound lower margins~8 and~7/8.
\paragraph{Conditions and connections.}
An affine objective's corner check can verify arithmetic only after stable
ReLU branches and affine dependence are established. The proof is tied to
the anchor's overlap; C24 shows why repeating the wave box at a new anchor
without new evidence can fail.

\concept{B29}{A Cartesian product region does not assert independence}{F24, F30, B12}
\paragraph{Foundation and definition.}
Geometry specifies permitted values; probability specifies how values are
distributed. A finite-horizon wave region may be the Cartesian product
\[
 \mathcal V=\mathcal U_1\times\cdots\times\mathcal U_Q.
\]
\paragraph{Symbols and meaning.}
$Q\ge1$ is the number of future steps, $\mathcal U_q\subseteq\R^{Cs}$
is the admissible wave set at step $q$, $C$ is channels and $s$ is
samples per wave. A member of $\mathcal V$ is a sequence
$(U_1,\ldots,U_Q)$ with $U_q\in\mathcal U_q$. No probability measure
is specified by this set notation.
\paragraph{Derivation and example.}
Let $Q=2$, $\mathcal U_1=\mathcal U_2=[-1,1]$, and random variables
satisfy $U_2=U_1$ with $U_1$ uniform on $[-1,1]$. Their pair always
belongs to the product square but lies only on its diagonal. The variables
are dependent, while a proof that holds throughout the square also holds
on that diagonal. Independence would be a factorization statement about
joint probabilities, not a consequence of the square's shape.
\paragraph{Why ECDR needs it.}
Finite-horizon verification can use a product region without assuming that
sensor waves are statistically independent. A dependent reachable set may
be smaller and easier to certify, but its constraints need a justified
membership or transition contract. A geometry that admits impossible
sequences can still be sound as an overapproximation, though less useful.
\paragraph{Conditions and connections.}
A product domain by itself supplies no coverage probability and no
guarantee that future waves remain inside it. Conversely, observed temporal
correlation is not a deterministic constraint on every future wave.
C25 combines transition geometry with induction; B12 handles joint events.

\concept{B30}{Classification metrics describe a specified labeled population}{F09, F31--F32, B06, B15}
\paragraph{Foundation and definition.}
Prediction quality against true labels differs from reuse agreement with
a reference. For a confusion matrix with true rows and predicted columns,
consider
\[
 M=\begin{pmatrix}8&2\\1&9\end{pmatrix}.
\]
\paragraph{Symbols and meaning.}
$M_{ij}$ counts examples truly in class $i$ and predicted as class $j$;
$N=\sum_{i,j}M_{ij}=20$ is population size. For a class, TP means true
positives, FP false positives and FN false negatives. Precision is
$\mathrm{TP}/(\mathrm{TP}+\mathrm{FP})$, recall is
$\mathrm{TP}/(\mathrm{TP}+\mathrm{FN})$, and
\[
 \mathrm{F1}=\frac{2\mathrm{TP}}{2\mathrm{TP}+\mathrm{FP}+\mathrm{FN}},
 \qquad \mathrm{MacroF1}=K^{-1}\sum_{j=1}^{K}\mathrm{F1}_j.
\]
$K$ is the number of classes included in that average, not an adjustable
denominator chosen after seeing results.
\paragraph{Derivation and example.}
Accuracy is $(8+9)/20=0.85$. Class1 has TP~8, FP~1, FN~2, so F1 is~16/19.
Class2 has F1 $18/21=6/7$. Macro-F1 is
$(16/19+6/7)/2=113/133\approx0.849624$, not exactly accuracy.
Zero denominators need an explicit convention. Across seeds, mean and
standard deviation are descriptive summaries; standard deviation alone
is not a confidence interval or an independent test split.
\paragraph{Why ECDR needs it.}
A reference may misclassify a true activity while a reuse mechanism
faithfully returns the same mistaken class. Accuracy, reference agreement,
certified coverage and net runtime benefit therefore answer different
questions. Filtering to pure K6 changes the labeled population relative
to full K12 or the all-window stream.
\paragraph{Conditions and connections.}
Declare labels, class schema, policy, split, support and aggregation before
interpreting a metric. Temporal dependence affects uncertainty estimates
even when the arithmetic is correct. C28 gives separate denominators for
certification and reuse.
